%=============================================================================
% Chapter 9: Database Design
%=============================================================================
\chapter{Database Design}

\section{Data Storage Strategy}

FloodEye employs a distributed data storage strategy consisting of three complementary layers rather than a single monolithic database:

\begin{enumerate}[leftmargin=1.5em]
  \item \textbf{ThingSpeak Cloud Storage:} The primary time-series store for all sensor readings. ThingSpeak acts as the IoT backend database, storing every field-mapped reading uploaded by the ESP32 with a UTC timestamp and an auto-incrementing \texttt{entry\_\allowbreak{}id}.
  \item \textbf{Browser localStorage:} A client-side persistence layer that stores the last known sensor reading, history array, device status, and last-seen timestamp. This enables offline resilience when the ESP32 or ThingSpeak is unreachable.
  \item \textbf{In-Memory Server Cache:} The Next.js server maintains an in-process cache (using \texttt{globalThis.\allowbreak{}\_\allowbreak{}\_\allowbreak{}thingspeakCache}) that prevents repeated API calls to ThingSpeak within the cache TTL window.
\end{enumerate}

\textit{Note: The project documentation (\texttt{TIHDash.\allowbreak{}md}) references a Neon PostgreSQL database for storing users, devices, alert history, and settings. In the current implementation, this database layer is not active. All session data, settings, and alerts are managed in-memory/localStorage only. A Neon PostgreSQL integration is identified as a future enhancement.}

\section{ThingSpeak Data Model}

\subsection{Channel Structure}

ThingSpeak organises data into channels, each containing up to eight fields. FloodEye uses five fields:

\begin{table}[H]
\centering
\caption{ThingSpeak Data Schema}
\begin{tabularx}{\linewidth}{|l|l|l|X|}
\hline
\rowcolor{FloodLight}
\textbf{Field} & \textbf{Name} & \textbf{Type} & \textbf{Description} \\
\hline
\texttt{field1} & temperature & FLOAT & DHT11 temperature reading in °C \\
\texttt{field2} & humidity & FLOAT & DHT11 relative humidity in \% \\
\texttt{field3} & pressure & FLOAT & BMP180 atmospheric pressure in hPa \\
\texttt{field4} & altitude & FLOAT & BMP180 derived altitude in metres \\
\texttt{field5} & distance & FLOAT & HC-SR04 distance to water surface in cm \\
\texttt{created\_\allowbreak{}at} & timestamp & DATETIME & UTC timestamp of the upload \\
\texttt{entry\_\allowbreak{}id} & id & INTEGER & Auto-increment entry identifier \\
\hline
\end{tabularx}
\end{table}

\subsection{Data Access Patterns}

Two primary query patterns are used:

\begin{enumerate}[leftmargin=1.5em]
  \item \textbf{Latest Entry:} \texttt{GET /\allowbreak{}channels/\allowbreak{}\{id\}/feeds/last.json?api\_key=\{key\}} --- Returns a single JSON object (the most recent feed entry). Used by \texttt{/\allowbreak{}api/\allowbreak{}sensors}.
  \item \textbf{Historical Range:}
    \begin{itemize}
      \item By count: \texttt{GET /\allowbreak{}channels/\allowbreak{}\{id\}/feeds.json?results=N\&api\_key=\{key\}} --- Returns the last N entries (max 8000 per call). Used for ranges $\leq$ 24h.
      \item By date range: \texttt{GET /\allowbreak{}channels/\allowbreak{}\{id\}/feeds.json?start=\{dt\}\&end=\{dt\}\&api\_key=\{key\}} --- Returns all entries in the specified UTC date window. Used for 7d and 30d ranges.
    \end{itemize}
\end{enumerate}

\section{Client-Side Data Model}

\subsection{TelemetryData Interface}

\begin{lstlisting}[language=TypeScript,caption={TelemetryData interface}]
export interface TelemetryData {
  temperature: number;  // deg C
  humidity:    number;  // %
  pressure:    number;  // hPa
  altitude:    number;  // metres
  distance:    number;  // cm (water distance)
  timestamp:   string;  // ISO 8601 UTC string
}
\end{lstlisting}

\subsection{DeviceStatus Interface}

\begin{lstlisting}[language=TypeScript,caption={DeviceStatus interface}]
export interface DeviceStatus {
  esp32Online:         boolean;
  thingspeakConnected: boolean;
  lastUpload:          string;  // ISO timestamp
  rssi:                number;  // Not available via REST (0)
}
\end{lstlisting}

\subsection{Alert Interface}

\begin{lstlisting}[language=TypeScript,caption={Alert interface}]
export interface Alert {
  id:        string;        // random alphanumeric
  type:      AlertType;     // HIGH_TEMPERATURE | HIGH_HUMIDITY | ...
  severity:  AlertSeverity; // "critical" | "warning" | "info"
  message:   string;
  timestamp: string;        // ISO creation time
  read:      boolean;
  updatedAt: string;        // ISO of last refresh
}
\end{lstlisting}

\subsection{DashboardSettings Interface}

\begin{lstlisting}[language=TypeScript,caption={DashboardSettings interface}]
export interface DashboardSettings {
  refreshInterval:   number;  // seconds (5 | 15 | 30 | 60)
  tempThreshold:     number;  // deg C  (default: 35)
  humidityThreshold: number;  // %       (default: 80)
  distanceThreshold: number;  // cm      (default: 20)
  pressureThreshold: number;  // hPa     (default: 1005)
}
\end{lstlisting}

\section{localStorage Schema}

\begin{table}[H]
\centering
\caption{localStorage Key Schema}
\begin{tabularx}{\linewidth}{|l|l|X|}
\hline
\rowcolor{FloodLight}
\textbf{Key} & \textbf{Value Type} & \textbf{Description} \\
\hline
\texttt{ecosense:lastData} & TelemetryData (JSON) & Most recent live reading \\
\texttt{ecosense:lastHistory} & TelemetryData[] (JSON) & Last known history array \\
\texttt{ecosense:lastStatus} & DeviceStatus (JSON) & Last known device status \\
\texttt{ecosense:lastSeenAt} & string (ISO) & Timestamp of last live contact \\
\texttt{floodeye\_\allowbreak{}session} & "admin" | "user" & Current user role \\
\texttt{floodeye\_\allowbreak{}user\_\allowbreak{}name} & string & Display name \\
\hline
\end{tabularx}
\end{table}

\section{ML Model Data Files}

The machine learning pipeline uses several JSON data files stored in \texttt{/\allowbreak{}public/\allowbreak{}} for browser access:

\begin{table}[H]
\centering
\caption{ML Data Files}
\begin{tabularx}{\linewidth}{|l|X|}
\hline
\rowcolor{FloodLight}
\textbf{File} & \textbf{Contents} \\
\hline
\texttt{public/\allowbreak{}model\_\allowbreak{}config.\allowbreak{}json} & LOOKBACK (48), DANGER\_THRESHOLD (450 cm), feature list \\
\texttt{public/\allowbreak{}model\_\allowbreak{}scalers.\allowbreak{}json} & StandardScaler means and scales for X (11 features) and Y (water level) \\
\texttt{public/\allowbreak{}tfjs\_\allowbreak{}model/\allowbreak{}model.\allowbreak{}json} & TensorFlow.js LayersModel topology and weights manifest \\
\texttt{public/\allowbreak{}tfjs\_\allowbreak{}model/\allowbreak{}group1-shard1of1.\allowbreak{}bin} & Float32 binary weights file (291 KB) \\
\texttt{public/\allowbreak{}sample\_\allowbreak{}data.\allowbreak{}json} & 200-row sample of test-set data used to pad insufficient history for inference \\
\hline
\end{tabularx}
\end{table}

\section{Data Retention and History}

ThingSpeak free-tier channels retain data for 1 year. The dashboard can query up to 30 days of history via the date-range API. The history panel displays the last 100 readings in the scrollable table. CSV/PDF export is limited to the last 100 readings in the current implementation.
